\documentclass[a4paper,12pt]{article}
\usepackage[utf8]{inputenc}
\usepackage{amsmath}
\usepackage{graphicx}
\usepackage{hyperref}
\usepackage{geometry}
\geometry{left=3cm, right=2.5cm, top=2.5cm, bottom=2.5cm}

\title{\textbf{Proposta de Projeto Final}}
\author{Nome do Estudante}
\date{\today}

\begin{document}

% Capa com imagem
\begin{titlepage}
    \centering
    \includegraphics[width=0.5\textwidth]{istec.jpeg}
    \vspace{1cm}

    {\LARGE \textbf{ISTEC} \par}
    \vspace{1.5cm}

    {\Large \textbf{Projeto Final} \par}
    \vspace{0.5cm}

    \textbf{Goncalo Queirós} \\
    Curso:Desenvolvimento de Software \\

    \textbf{Disciplina} \\
    Desenvolvimento Agil de Software \\

    \vfill

    {\large \today\par}
\end{titlepage}

\newpage
\tableofcontents
\newpage

\section{Introdução}
Este relatório descreve o processo de criação de um repositório GitHub utilizando o modelo GitFlow para gestão de branches, controle de versões e revisão de código. Foram aplicadas práticas recomendadas para gerenciar acessos, configurar revisões obrigatórias e ignorar arquivos desnecessários.

\section{Criação do repositório e configuração do gitflow}


\fbox{
\begin{minipage}{1\linewidth}

c c:repositorios\\
cd projetofinal\\

git init\\

git remote add origin http:/etc\\

git add .\\
git commit -m "primeiro commit"\\

git branch -m main\\
git push -u origin main\\

git flow init\\



git push origin master\\

\end{minipage}
}



\section{Desenvolvimento}
 \fbox{
\begin{minipage}{1\linewidth}

git flow feature start feature1\\
git flow feature finish feature1\\

git flow hotfix start versao1\\
git flow hotfix finish versao1\\

git flow bugfix start versao1\\
git flow bugfix finish versao1\\

git add \texttt{.}gitignore\\
echo "*.doc" \texttt{>>} .gitignore\\
echo "*.docx" \texttt{>>} .gitignore\\

\end{minipage}
}


\end{document}